\documentclass[12pt,a4paper]{article}

\usepackage[T1]{fontenc}
\usepackage[utf8]{inputenc}
\usepackage[brazil]{babel}
\usepackage{lmodern}
\usepackage{geometry}
\usepackage{hyperref}
\usepackage{enumitem}
\usepackage{array}

\geometry{margin=2.5cm}
\hypersetup{
  colorlinks=true,
  linkcolor=blue,
  urlcolor=blue
}

\title{Plano de Aula: Operações da Álgebra Relacional}
\date{}

\begin{document}
\maketitle

\noindent
\textbf{Duração:} 2 horas (120 minutos)\\
\textbf{Público-alvo:} Iniciantes\\
\textbf{Recursos Necessários:} Quadro branco/negro, marcadores/giz, cartões/papel, post-its, fita adesiva e canetas.

\section{Objetivos da Aula}
\begin{itemize}[leftmargin=*,itemsep=0.25em]
  \item Apresentar as principais operações da álgebra relacional.
  \item Interpretar consultas como operações sobre conjuntos de tuplas.
  \item Relacionar álgebra relacional com a ideia de consultas em SQL.
  \item Resolver problemas simples usando seleção, projeção e junção (em nível conceitual).
\end{itemize}

\section{Metodologia}
Aula conceitual com simulações físicas (tabelas no quadro e cartões como tuplas). Os alunos ``executam'' operações movendo/filtrando cartões, criando o entendimento de consulta antes da sintaxe do SQL.

\section{Cronograma e Conteúdo Detalhado (120 Minutos)}
\subsection*{Parte 1: Introdução --- Consultas como ``operações'' (15 minutos)}
\begin{itemize}[leftmargin=*,itemsep=0.35em]
  \item Relembrar: tabela como conjunto de tuplas.
  \item Analogia: peneirar (filtrar), recortar colunas (selecionar atributos), juntar tabelas (conectar informações).
\end{itemize}

\subsection*{Parte 2: Operações Fundamentais (35 minutos)}
\begin{itemize}[leftmargin=*,itemsep=0.35em]
  \item \textbf{Seleção ($\sigma$) (10 min):} filtra linhas por condição.
  \item \textbf{Projeção ($\pi$) (10 min):} seleciona colunas (remove atributos).
  \item \textbf{Produto cartesiano e junção (10 min):} combinar linhas de duas relações.
  \item \textbf{União/Diferença (5 min):} operações entre conjuntos compatíveis.
\end{itemize}

\subsection*{Parte 3: Exemplo Guiado no Quadro (20 minutos)}
\begin{itemize}[leftmargin=*,itemsep=0.35em]
  \item Usar tabelas CLIENTE e PET.
  \item Mostrar uma seleção: ``somente clientes da cidade X''.
  \item Mostrar uma projeção: ``apenas nomes dos clientes''.
  \item Mostrar junção conceitual: ``pets com nome do cliente''.
\end{itemize}

\subsection*{Parte 4: Atividade em Grupo --- Álgebra com Cartões (40 minutos)}
\begin{itemize}[leftmargin=*,itemsep=0.35em]
  \item \textbf{Preparação:} distribuir cartões representando tuplas de duas tabelas (ex.: CLIENTE e PET).
  \item \textbf{Tarefas (cada grupo executa e registra a resposta):}
  \begin{itemize}[leftmargin=*,itemsep=0.25em]
    \item Aplicar $\sigma$ para filtrar registros por condição definida pelo professor.
    \item Aplicar $\pi$ para retornar apenas colunas específicas.
    \item Fazer uma junção ``manual'' usando a FK para combinar informações.
  \end{itemize}
  \item \textbf{Entrega:} resultado das operações (quais tuplas aparecem).
\end{itemize}

\subsection*{Parte 5: Encerramento e Ponte para SQL (10 minutos)}
\begin{itemize}[leftmargin=*,itemsep=0.35em]
  \item ``SQL é uma forma de pedir ao SGBD para executar essas operações''.
  \item Ponte: na próxima aula, sintaxe básica de SQL para criar e manipular dados.
\end{itemize}

\end{document}

